\documentclass[10pt,a4paper]{amsart}
\usepackage[margin=19mm]{geometry}
\usepackage{amsmath,amssymb,mathtools}
\usepackage[scr=rsfs,cal=euler]{mathalfa}
\usepackage{xfrac}
\usepackage[unicode,hidelinks,bookmarks=false]{hyperref}
\newcommand{\ie}{i.e.\ }
\newcommand{\cf}{cf.\ }
\newcommand{\resp}{respectively\ }
\newcommand{\Subsection}{\subsection}
\providecommand{\nobreakdash}{\nobreak\hbox{-}}
\setlength{\emergencystretch}{2em}
\multlinegap0pt
\usepackage{mathrsfs,colonequals}
\hypersetup{colorlinks=true,urlcolor=magenta,citecolor=magenta,linkcolor=blue}
\newcommand{\Gal}{\operatorname{Gal}}
\newcommand{\tr}{\operatorname{tr}}
\newcommand{\Z}{\mathbf{Z}}
\newcommand{\R}{\operatorname{\mathbf{R}}}
\newcommand{\Q}{\mathbf{Q}}
\newcommand{\ddef}{\colonequals}
\newcommand{\disc}{\operatorname{disc}}
\newtheorem{thm}[equation]{Theorem}
\newtheorem{lem}[equation]{Lemma}
\newtheorem{defn}[equation]{Definition}
\newtheorem{cor}[equation]{Corollary}
\newtheorem{prop}[equation]{Proposition}
\newtheorem{rmk}[equation]{Remark}
\newtheorem{remark}[equation]{Remark}
\newtheorem{exm}[equation]{Example}
\begin{document}
\title[Hasse-Witt invariants for trace forms of Jacobi polynomials]{Hasse-Witt invariants for trace forms of Jacobi polynomials}
\author{John Cullinan, Farshid Hajir, Elisabeth Young}
\date{}
\maketitle
\noindent\emph{Source and licence.} Hasse-Witt invariants for trace forms of Jacobi polynomials. Version arXiv:2608.12019v1 (CC BY 4.0). This material is distributed under the Creative Commons Attribution 4.0 International licence, http://creativecommons.org/licenses/by/4.0/, as recorded in the arXiv OAI licence metadata for this exact version and shown on the abstract page https://arxiv.org/abs/2608.12019v1. Full rights notice: licenses/arxiv-2608.12019-CC-BY-4.0.txt.\par\smallskip
\noindent\textit{Editorial scope.} Counted unit: the original \texttt{thm} environment \texttt{elis\_version} at source lines 411--423 together with its complete proof at lines 425--449; the delivered document keeps source lines 151--466 so that every referenced label and every citation resolves inside it. Native editable source is the paper's own concatenated TeX; the AMS wrapper is editorial formatting only. No publisher class, upstream script or unrelated artwork is redistributed.
\begin{align}
A_{t} &= \prod_{j=n-t+1}^n j^{j-(n-t)}, \label{ADEF} \\
B_{t} &= \prod_{j=n-t+2}^n \left[(\alpha + j)(\beta+j)\right]^{j-(n-t+1)} \left(\alpha + \beta + j\right)^{j-(n-t+2)}, \label{BDEF} \\
C_{t} &= \prod_{j=2n-2t+2}^{2n} (\alpha + \beta + j)^{j-(2n-2t+2)}. \label{CDEF}
\end{align}
With this notation in place, we can now state the main results of this paper.

\begin{thm} \label{main_thm}
Fix $n \geq 2$.  For the polynomial  $\mathscr{J}_n^{(\alpha,\beta)}(x)$, we have $\Delta_{t} = A_{t}B_{t}/C_{t}$, with $A_t$, $B_t$, and $C_t$  as defined in Equations (\ref{ADEF}), (\ref{BDEF}), and (\ref{CDEF}).
\end{thm}

In \cite[Thm.~3.1]{feit}, Feit proved the general formula
\begin{align} \label{HW_general}
\epsilon_p(E) = \left\{ \prod_{t=1}^n \left( \Delta_{t-1},\Delta_{t} \right)_p \right\} \left( -1,\prod_{t=1}^n \Delta_{t} \right)_p,
\end{align}
where $(~,~)_p$ denotes the $p$-adic Hilbert symbol.  By applying Theorem \ref{main_thm} to Equation (\ref{HW_general}) we will obtain in Section \ref{HW_section} an explicit expression for the Hasse-Witt invariant of the root field of $\mathscr{J}_n^{(\alpha,\beta)}(x)$.  We then show how to recover \cite[Thm.~5.1]{feit} immediately from Theorem \ref{main_thm} and give several concrete examples pertaining to the Jacobi Polynomials.   

\section{Hankel Matrices}

\subsection{Determinants} A \textbf{Hankel matrix} is one that is constant along anti-diagonals:
\[
\begin{pmatrix} 
a & b & c & \cdots \\
b & c & d & \cdots \\
c & d & e &  \cdots \\
\vdots & \vdots & \vdots & \ddots
\end{pmatrix}.
\]
If $\lbrace \mu_k \rbrace_{k\geq 0}$ is a sequence, then we can define the associated $t \times t$ Hankel matrix via 
\[
\left( \mu_{i+j} \right)_{0 \leq i,j \leq t-1}.
\]
Given such a Hankel matrix, one is led to consider its determinant 
\begin{align} \label{general_det}
\det \left(\mu_{i+j}\right)_{0 \leq i,j \leq t-1}.
\end{align} 
In certain circumstances, if the sequence $\lbrace \mu_k \rbrace_{k \geq 0}$ is endowed with some extra structure, there are techniques for evaluating the determinant in (\ref{general_det}). We recall one of those circumstances now.

Let $\lbrace \mu_k \rbrace_{k \geq 0}$ be a sequence with generating function
\[
G(x) =  \sum_{k=0}^\infty \mu_kx^k.
\]
If we write $G(x)$ as a continued fraction of the form
\begin{align} \label{cont_frac}
G(x) = \sum_{k=0}^\infty \mu_kx^k  = \frac{\mu_0}{\displaystyle 1+  \frac{a_1x}{1+ \displaystyle \frac{a_2x}{1 + \cdots}}},
\end{align}
then by \cite[Thm.~30]{krattenhaler}, we have
\begin{align} \label{kratt_form}
\det \left( \mu_{i+j} \right)_{0 \leq i,j \leq t-1} = \mu_0^t \left( a_1a_2 \right)^{t-1}\left(a_3a_4 \right)^{t-2} \cdots \left(a_{2t-3}a_{2t-2}\right).
\end{align}
The minor $\Delta_{t}$ introduced above is a Hankel determinant, hence one approach to proving Theorem \ref{main_thm} is to determine the continued fraction expansion as in Equation (\ref{cont_frac}) for the generating function of the power-sums.

\begin{rmk}
In \cite{feit}, Feit considered the determinant 
\[
\det \left(s_{i+j-2} \right)_{1 \leq i,j \leq t}.
\]
In order to have Feit's indexing agree with the indexing above, we set $i \mapsto i-1$ and $j \mapsto j-1$ and start the indexing at 0.
\end{rmk}

\subsection{Generating Functions}

Let $F \in \Q[x]$ be a polynomial of degree $n \geq 1$ with roots $\theta_1,\dots,\theta_n \in \overline{\Q}$.  Let $G(z) = \sum_{i=0}^\infty s_iz^i$ be the generating function for the power-sums of the roots.  For any polynomial $f$, let $\widehat{f}$ denote the reciprocal polynomial.

\begin{prop} \label{G_deriv_prop}
With all notation as above, we have
\[
G(z) = \frac{\widehat{F'(z)}}{\widehat{F(z)}}.
\]
\end{prop}

\begin{proof}
Fix an index $j \in \lbrace 1,\dots,n \rbrace$.  Then
\[
\frac{1}{1-z\theta_j} = \sum_{i=0}^\infty \theta_j^iz^i.
\]
Therefore, 
\[
G(z) = \sum_{i=0}^\infty s_iz^i = \sum_{i=0}^\infty \sum_{j=1}^n \theta_j^i z^i = \sum_{j=1}^n \sum_{i=0}^\infty \theta_j^iz^i= \sum_{j=1}^n \frac{1}{1-z\theta_j} =  \frac{\widehat{F'(z)}}{\widehat{F(z)}}.
\]
\end{proof}

Our strategy for determining $\Delta_t$ will be to write $G(z) = {\widehat{F'(z)}}/{\widehat{F(z)}}$ in the form of Equation (\ref{cont_frac}), thereby determining the coefficients $a_i$. Then
\begin{align} \label{det_k}
\Delta_{t} = s_0^t \left( a_1a_2 \right)^{t-1}\left(a_3a_4 \right)^{t-2} \cdots \left(a_{2t-3}a_{2t-2}\right),
\end{align}
with $s_0=n$.

\section{Proof of Theorem \ref{main_thm}}

For ease of exposition, let 
\begin{align} \label{gamma_defn}
\gamma_j(n)^{(\alpha,\beta)} = \binom{n}{j} \left( \prod_{k=j+1}^n \frac{\alpha + k}{n+\alpha + \beta + k} \right) (-1)^{n-j},
\end{align}
so that $\mathscr{J}_n^{(\alpha,\beta)}(x) = \sum_{j=0}^n\gamma_j(n)^{(\alpha,\beta)}x^j$. We start by writing $G(z)$ in the form of Proposition \ref{G_deriv_prop}.  


\begin{prop}\label{gen_function}
Let $n \geq 1$.  The generating function $G(z)$ for the power-sum polynomials in the roots of $\mathscr{J}_n^{(\alpha,\beta)}(x)$ is given by
\begin{align}  \label{F_eq_original}
G(z) =\frac{\sum_{j=0}^{n-1} \gamma_j^{(\alpha+1,\beta+1)}(n-1)z^{n-1-j}}{\sum_{j=0}^n \gamma_j^{(\alpha,\beta)}(n)z^{n-j}}.
\end{align}
\end{prop}

\begin{proof}
Fix $n \geq 1$.  By Proposition \ref{G_deriv_prop}, 
\[
G(z) = \frac{ \widehat{ \mathscr{J}_n^{(\alpha,\beta)}(z)'}}{\widehat{ \mathscr{J}_n^{(\alpha,\beta)}(x)}}.
\]
It is well known that
\[
\frac{{\rm d}}{{\rm d}z} P_n^{(\alpha,\beta)}(z) = \frac{\Gamma(\alpha+\beta+n+2)}{2\Gamma(\alpha+\beta+n+1)} P_{n-1}^{(\alpha+1,\beta+1)}(z),
\]
which gives us
\[
\frac{{\rm d}}{{\rm d}z} \mathscr{J}_n^{(\alpha,\beta)}(z) = n \mathscr{J}_{n-1}^{(\alpha + 1,\beta+1)}(z).
\]
Equipped with this derivative formula, the claimed expression for $G(z)$ follows by direct calculation. 
\end{proof}

We now turn to the proof of Theorem \ref{main_thm}.  To do this we will determine the $a_i$ of Equation (\ref{det_k}); it will then follow by algebraic manipulation that Equation (\ref{det_k}) and the one of Theorem \ref{main_thm} are the same.

Fix $n \geq 1$ and set the following notation. Let $c_i=i\gamma_i^{(\alpha,\beta)}(n)$ and $d_i=\gamma_i^{(\alpha,\beta)}(n)$ for all $0\le i\le n$, so that $c_i = id_i$. We define $f_0,\ldots,f_{2n-1}:\{0,\dots,n\} \to \R$ as follows. Let 
\begin{align}
f_0(m)&=\begin{cases} c_m,&1\le m\le n\\
0,&m=0, \end{cases} \label{f_0_eq}
\\
f_1(m)&=\begin{cases} d_{m-1}-\frac{f_0(m-1)}{f_0(n)},&1\le m\le n\\
0,&m=0,\hspace{0.5in}\text{and} \end{cases} \label{f_1_eq}
\\
f_i(m)&=\begin{cases}
\frac{f_{i-2}(m-1)}{f_{i-2}(n)}-\frac{f_{i-1}(m-1)}{f_{i-1}(n)},&\frac{i}{2}< m\le n\\
0,&0\le m\le\frac{i}{2}
\end{cases} \label{f_i_eq} \end{align}
for $2\le i\le 2n-1$. Let $a_i=f_i(n)$ for  $0\le i\le 2n-1$. We also note that for  $1\le m\le n$ we can simplify Equation (\ref{f_1_eq}) to get
\begin{align} \label{f_1_simplified}
    f_1(m)&=\gamma_{m-1}^{(\alpha,\beta)}(n)-\frac{(m-1)\gamma_{m-1}^{(\alpha,\beta)}(n)}{n} =\gamma_{m-1}^{(\alpha,\beta)}(n)\cdot\frac{n-m+1}{n}.
\end{align}

\begin{lem} \label{F_eq_cd}
We have \[G(z)=\frac{\sum_{j=1}^{n}c_jz^{n-j}}{\sum_{j=0}^nd_jz^{n-j}}.\]
\end{lem}

\begin{proof}
Using Equation (\ref{gamma_defn}) we have
\begin{align*}
    \gamma_j^{(\alpha+1,\beta+1)}(n-1)=\frac{j+1}{n}\gamma_{j+1}^{(\alpha,\beta)}(n),
\end{align*}
which we substitute into Equation (\ref{F_eq_original}) to get
\begin{align*}
    G(z)=\frac{\sum_{j=1}^{n}j\gamma_j^{(\alpha,\beta)}(n)z^{n-j}}{\sum_{j=0}^n\gamma_j^{(\alpha,\beta)}(n)z^{n-j}}.
\end{align*}
Applying the definitions of $c_i$ and $d_i$ completes the proof.
\end{proof}

\begin{prop} \label{cont_frac_prop}
The generating function $G(z)$ of Proposition \ref{gen_function} is given by  \[G(z)=\dfrac{a_0}{1+\dfrac{a_1z}{1+\dfrac{a_2z}{\dfrac{\vdots}{1+\dfrac{a_{2n-2}z}{1+a_{2n-1}z}}}}},\]
with $a_i=f_i(n)$ defined as above.
\end{prop}

\begin{proof}
We will show that 
\begin{align*}
G(z)=\dfrac{a_0}{1+\dfrac{a_1z}{1+\dfrac{a_2z}{\dfrac{\vdots}{1+\dfrac{a_{2q-1}z}{\dfrac{\sum_{j=q}^n\frac{f_{2q-2}(j)}{f_{2q-2}(j)}z^{n-j}}{\sum_{j=q}^n\frac{f_{2q-1}(j)}{f_{2q-1}(n)}z^{n-j}}}}}}}
\end{align*}
for all $1\le q\le n$; setting $q=n$ will complete the proof. We will induct on $q$. 



By Lemma \ref{F_eq_cd} we have
\begin{align*}
    G(z)=\frac{\sum_{j=1}^{n}c_jz^{n-j}}{\sum_{j=0}^nd_jz^{n-j}}.
\end{align*}
Note that $c_j=f_0(j)$ for all $1\le j\le n$. We also have $c_n=f_0(n)=a_0$ and $d_n=\gamma_n^{(\alpha,\beta)}(n)=1$. In particular, $a_0 = n$. Hence
\begin{align*}
    G(z)&=\frac{\bigg(\sum_{j=1}^{n-1}f_0(j)z^{n-j}\bigg)+a_0}{\bigg(\sum_{j=0}^{n-1}d_jz^{n-j}\bigg)+1} =\dfrac{a_0}{\dfrac{\bigg(\sum_{j=0}^{n-1}d_jz^{n-j}\bigg)+1}{\bigg(\sum_{j=1}^{n-1}\frac{f_0(j)}{f_0(n)}z^{n-j}\bigg)+1}}.
\end{align*}
Observe that 
\begin{align*}
{\dfrac{\bigg(\sum_{j=0}^{n-1}d_jz^{n-j}\bigg)+1}{\bigg(\sum_{j=1}^{n-1}\frac{f_0(j)}{f_0(n)}z^{n-j}\bigg)+1}} &= {1+\dfrac{\bigg(\sum_{j=0}^{n-1}(d_j-\frac{f_0(j)}{f_0(n)})z^{n-j}\bigg)}{\bigg(\sum_{j=1}^{n-1}\frac{f_0(j)}{f_0(n)}z^{n-j}\bigg)+1}} \\
&={1+\dfrac{z\bigg(\sum_{j=1}^nf_1(j)z^{n-j}\bigg)}{\bigg(\sum_{j=1}^{n-1}\frac{f_0(j)}{f_0(n)}z^{n-j}\bigg)+1}}.
\end{align*}
Thus,
\begin{align*}    
     G(z)&=\dfrac{a_0}{1+\dfrac{z\bigg(\sum_{j=1}^nf_1(j)z^{n-j}\bigg)}{\bigg(\sum_{j=1}^{n-1}\frac{f_0(j)}{f_0(n)}z^{n-j}\bigg)+1}}
    =\dfrac{a_0}{1+\dfrac{a_1z}{\dfrac{\sum_{j=1}^n\frac{f_0(j)}{f_0(n)}z^{n-j}}{\sum_{j=1}^n\frac{f_1(j)}{f_1(n)}z^{n-j}}}},
\end{align*}
which completes the base case. Fix $r\in\{1,\ldots,n-1\}$. Suppose 
\begin{align*}
G(z)=\dfrac{a_0}{1+\dfrac{a_1z}{1+\dfrac{a_2z}{\dfrac{\vdots}{1+\dfrac{a_{2r-1}z}{\dfrac{\sum_{j=r}^n\frac{f_{2r-2}(j)}{f_{2r-2}(n)}z^{n-j}}{\sum_{j=r}^n\frac{f_{2r-1}(j)}{f_{2r-1}(n)}z^{n-j}}}}}}}.
\end{align*}
Then
\begin{align*}
{1+\dfrac{a_{2r-1}z}{\dfrac{\sum_{j=r}^n\frac{f_{2r-2}(j)}{f_{2r-2}(n)}z^{n-j}}{\sum_{j=r}^n\frac{f_{2r-1}(j)}{f_{2r-1}(n)}z^{n-j}}}} &=
{1+\dfrac{a_{2r-1}z}{1+\dfrac{\sum_{j=r}^{n-1}\bigg(\frac{f_{2r-2}(j)}{f_{2r-2}(n)}-\frac{f_{2r-1}(j)}{f_{2r-1}(n)}\bigg)z^{n-j}}{\sum_{j=r}^n\frac{f_{2r-1}(j)}{f_{2r-1}(n)}z^{n-j}}}} \\
&={1+\dfrac{a_{2r-1}z}{1+\dfrac{a_{2r}z}{\dfrac{\sum_{j=r}^n\frac{f_{2r-1}(j)}{f_{2r-1}(n)}z^{n-j}}{\sum_{j=r+1}^{n}\frac{f_{2r}(j)}{f_{2r}(n)}z^{n-j}}}}} \\
&={1+\dfrac{a_{2r-1}z}{1+\dfrac{a_{2r}z}{1+\dfrac{\sum_{j=r}^{n-1}\big(\frac{f_{2r-1}(j)}{f_{2r-1}(n)}-\frac{f_{2r}(j)}{f_{2r}(n)}\big)z^{n-j}}{\sum_{j=r+1}^{n}\frac{f_{2r}(j)}{f_{2r}(n)}z^{n-j}}}}}\\
&={1+\dfrac{a_{2r-1}z}{1+\dfrac{a_{2r}z}{1+\dfrac{a_{2r+1}z}{\dfrac{\sum_{j=r+1}^{n}\frac{f_{2r}(j)}{f_{2r}(n)}z^{n-j}}{\sum_{j=r+1}^{n}\frac{f_{2r+1}(j)}{f_{2r+1}(n)}z^{n-j}}}}}},
\end{align*}
which completes the proof.
\end{proof}

Next we determine alternate expressions for the $f_i$ that are more convenient for explicit calculation.

\begin{prop} \label{f(n-l)_prop}
For all $0\le i\le n-1$ and all $1\le \ell<n-\frac{i}{2}$, we have
\begin{align*}
    f_{2i}(n-\ell)&=f_{2i}(n-\ell+1)\cdot\frac{-(n-i-\ell)(\alpha+n-i-\ell+1)}{\ell(\alpha+\beta+2n-2i-\ell+1)}, \text{ and } \\
    f_{2i+1}(n-\ell)&=f_{2i+1}(n-\ell+1)\cdot\frac{-(n-i-\ell)(\alpha+n-i-\ell)}{\ell(\alpha+\beta+2n-2i-\ell)}.
\end{align*}    
\end{prop}

\begin{proof}
We will proceed by induction on $i$, beginning with the cases $i \in \lbrace 0,1\rbrace$ separately.  Applying  the definition of the $\gamma_{m}^{(\alpha,\beta)}$ we have
\begin{align*}
    f_0(n-\ell)&=(n-\ell)\gamma_{n-\ell}^{(\alpha,\beta)}(n)\\
    &=(n-\ell)\binom{n}{n-\ell}\bigg(\prod_{k=n-\ell+1}^n\frac{\alpha+k}{\alpha+\beta+n+k}\bigg)(-1)^\ell \\
    &=\frac{-(n-\ell)(n-\ell+1)}{\ell}\binom{n}{n-\ell+1}\bigg(\prod_{k=n-\ell+2}^n\frac{\alpha+k}{\alpha+\beta+n+k}\bigg)\bigg(\frac{\alpha+n-\ell+1}{\alpha+\beta+2n-\ell+1}\bigg)(-1)^{\ell-1}\\
    &=(n-\ell+1)\gamma_{n-\ell+1}^{(\alpha,\beta)}(n)\cdot\frac{-(n-\ell)(\alpha+n-\ell+1)}{\ell(\alpha+\beta+2n-\ell+1)}\\
    &=f_0(n-\ell+1)\cdot\frac{-(n-\ell)(\alpha+n-\ell+1)}{\ell(\alpha+\beta+2n-\ell+1)}, \text{ and} \\
    f_1(n-\ell)&=\gamma_{n-\ell-1}^{(\alpha,\beta)}(n)\frac{\ell+1}{n} \qquad \text{(by Equation (\ref{f_1_simplified}))}\\
    &=\binom{n}{n-\ell-1}\bigg(\prod_{k=n-\ell}^n\frac{\alpha+k}{\alpha+\beta+n+k}\bigg)(-1)^{\ell+1}\cdot\frac{\ell+1}{n}\\
    &=\binom{n}{n-\ell}\bigg(\prod_{k=n-\ell+1}^n\frac{\alpha+k}{\alpha+\beta+n+k}\bigg)(-1)^{\ell}\cdot\frac{-(n-\ell)(\alpha+n-\ell)}{n(\alpha+\beta+2n-\ell)}\\
    &=\gamma_{n-\ell}^{(\alpha,\beta)}(n)\cdot\frac{\ell}{n}\cdot\frac{-(n-\ell)(\alpha+n-\ell)}{\ell(\alpha+\beta+2n-\ell)}\\
    &=f_1(n-\ell+1)\cdot\frac{-(n-\ell)(\alpha+n-\ell)}{\ell(\alpha+\beta+2n-\ell)},
\end{align*}
which completes the base cases. Now fix $s\in\{0,\ldots,n-2\}$ and suppose 
\begin{align} \label{f_2s_induct_hyp}
    f_{2s}(n-\ell)&=f_{2s}(n-\ell+1)\cdot\frac{-(n-s-\ell)(\alpha+n-s-\ell+1)}{\ell(\alpha+\beta+2n-2s-\ell+1)},\text{ and} \\
 \label{f_2s+1_induct_hyp}
    f_{2s+1}(n-\ell)&=f_{2s+1}(n-\ell+1)\cdot\frac{-(n-s-\ell)(\alpha+n-s-\ell)}{\ell(\alpha+\beta+2n-2s-\ell)},
\end{align}
for all $1\le \ell<n-\frac{s}{2}$. Using Equations (\ref{f_i_eq}), (\ref{f_2s_induct_hyp}), and (\ref{f_2s+1_induct_hyp}) we have
\begin{align}
f_{2s+2}(n-\ell)&=\frac{f_{2s}(n-\ell-1)}{f_{2s}(n)}-\frac{f_{2s+1}(n-\ell-1)}{f_{2s+1}(n)}\notag \\
&=\bigg(\prod_{k=1}^{\ell+1}\frac{-(n-s-k)(\alpha+n-s-k+1)}{k(\alpha+\beta+2n-2s-k+1)}\bigg)-\bigg(\prod_{k=1}^{\ell+1}\frac{-(n-s-k)(\alpha+n-s-k)}{k(\alpha+\beta+2n-2s-k)}\bigg)\notag \\
&=\bigg(\prod_{k=1}^{\ell+1}\frac{-(n-s-k)}{k}\bigg)\bigg(\prod_{k=1}^\ell\frac{\alpha+n-s-k}{\alpha+\beta+2n-2s-k}\bigg) \notag 
\\&\hspace{1in}\cdot\bigg(\frac{\alpha+n-s}{\alpha+\beta+2n-2s}-\frac{\alpha+n-s-\ell-1}{\alpha+\beta+2n-2s-\ell-1}\bigg)\notag  \\
&=\bigg(\prod_{k=1}^{\ell+1}\frac{-(n-s-k)}{k}\bigg)\bigg(\prod_{k=1}^\ell\frac{\alpha+n-s-k}{\alpha+\beta+2n-2s-k}\bigg) \notag 
\\&\hspace{1in}\cdot\bigg(\frac{(\ell+1)(\beta+n-s)}{(\alpha+\beta+2n-2s)(\alpha+\beta+2n-2s-\ell-1)}\bigg)\notag \\
&=\bigg(\prod_{k=1}^{\ell}\frac{-(n-s-k)}{k}\bigg)\bigg(\prod_{k=1}^{\ell-1}\frac{\alpha+n-s-k}{\alpha+\beta+2n-2s-k}\bigg) \notag 
\\&\hspace{0.5in}\cdot\bigg(\frac{-(\beta+n-s)(n-s-\ell-1)(\alpha+n-s-\ell)}{(\alpha+\beta+2n-2s)(\alpha+\beta+2n-2s-\ell-1)(\alpha+\beta+2n-2s-\ell)}\bigg)\notag \\
&=\bigg(\prod_{k=1}^{\ell}\frac{-(n-s-k)}{k}\bigg)\bigg(\prod_{k=1}^{\ell-1}\frac{\alpha+n-s-k}{\alpha+\beta+2n-2s-k}\bigg) \notag 
\\&\hspace{0.5in}\cdot\bigg(\frac{\alpha+n-s}{\alpha+\beta+2n-2s}-\frac{\alpha+n-s-\ell}{\alpha+\beta+2n-2s-\ell}\bigg)\bigg(\frac{-(n-s-\ell-1)(\alpha+n-s-\ell)}{\ell(\alpha+\beta+2n-2s-\ell-1)}\bigg)\notag \\
&=\bigg[\bigg(\prod_{k=1}^\ell\frac{-(n-s-k)(\alpha+n-s-k+1)}{k(\alpha+\beta+2n-2s-k+1)}\bigg)-\bigg(\prod_{k=1}^\ell\frac{-(n-s-k)(\alpha+n-s-k)}{k(\alpha+\beta+2n-2s-k)}\bigg)\bigg] \notag 
\\&\hspace{0.5in}\cdot\bigg(\frac{-(n-s-\ell-1)(\alpha+n-s-\ell)}{\ell(\alpha+\beta+2n-2s-\ell-1)}\bigg)\notag \\
&=\bigg(\frac{f_{2s}(n-\ell)}{f_{2s}(n)}-\frac{f_{2s+1}(n-\ell)}{f_{2s+1}(n)}\bigg)\bigg(\frac{-(n-s-\ell-1)(\alpha+n-s-\ell)}{\ell(\alpha+\beta+2n-2s-\ell-1)}\bigg). \notag
\end{align}
This shows that 
\begin{align}
f_{2s+2}(n-\ell)&=f_{2s+2}(n-\ell+1)\cdot\frac{-(n-s-\ell-1)(\alpha+n-s-\ell)}{\ell(\alpha+\beta+2n-2s-\ell-1)}. \label{f_2s+2_induct_step}
\end{align}
A similar argument using Equations (\ref{f_i_eq}), (\ref{f_2s+1_induct_hyp}), and (\ref{f_2s+2_induct_step}) shows that
\begin{align}
f_{2s+3}(n-\ell)
&=f_{2s+3}(n-\ell+1)\cdot\frac{-(n-s-\ell-1)(\alpha+n-s-\ell-1)}{\ell(\alpha+\beta+2n-2s-\ell-2)}. \label{f_2s+3_induct_step}
\end{align}
This completes the proof.
\end{proof}


\begin{thm} \label{elis_version}
Fix $n \geq 1$.  Let $\{s_k\}_{k\geq 0}$ be the sequence of power-sums for the roots of  $\mathscr{J}_n^{(\alpha,\beta)}(x)$.  Then 
\begin{gather*}
    \det(s_{i+j})_{0\le i,j\le t-1}={s_0}^t(a_1a_2)^{t-1}
   (a_3a_4)^{t-2}\cdots(a_{2n-3}a_{2n-2}),
\end{gather*}
where $s_0 = a_0=n$, ${a_1=\frac{-(\alpha+n)}{(\alpha+\beta+2n)}}$, 
\begin{align*}
    a_{2i}&=\frac{-(n-i)(\beta+n-i+1)}{(\alpha+\beta+2n-2i+2)(\alpha+\beta+2n-2i+1)}, \\
    a_{2i+1}&=\frac{-(\alpha+n-i)(\alpha+\beta+n-i+1)}{(\alpha+\beta+2n-2i+1)(\alpha+\beta+2n-2i)},\text{ and}\\
\end{align*}
for  $1\le i\le n-1$.
\end{thm}

\begin{proof}
Recall from Equations (\ref{cont_frac}) and (\ref{kratt_form}) that given the generating function $G(z)$ for any sequence $\lbrace \mu_k \rbrace_{k \geq 0}$ written as a continued fraction, the associated Hankel determinant is given by
\begin{align} \label{det_form_defn}
    \det(\mu_{i+j})_{0\le i,j\le t-1}={\mu_0}^t(a_1a_2)^{t-1}
   (a_3a_4)^{t-2}\cdots(a_{2n-3}a_{2n-2}).
\end{align}
Having proved Proposition \ref{cont_frac_prop}, it remains to provide explicit expressions for the coefficients $a_0,\ldots a_{2n-2}$.  We have already observed that $s_0 = a_0 = n$.  By  Equation (\ref{f_1_simplified}) we have
\begin{align*}
    a_1=f_1(n)=\frac{\gamma_{n-1}^{(\alpha,\beta)}(n)}{n}=\frac{-(\alpha+n)}{\alpha+\beta+2n}.
\end{align*}
Fix $p\in\{1,\ldots,n-1\}$. Using Equation (\ref{f_i_eq}) and Proposition \ref{f(n-l)_prop} we have
\begin{align*}
    a_{2p}&=f_{2p}(n)\\
    &=\frac{f_{2p-2}(n-1)}{f_{2p-2}(n)}-\frac{f_{2p-1}(n-1)}{f_{2p-1}(n)}\\
    &=\frac{-(n-p)(\alpha+n-p+1)}{\alpha+\beta+2n-2p+2}-\frac{-(n-p)(\alpha+n-p)}{\alpha+\beta+2n-2p+1}\\
    &=\frac{-(n-p)(\beta+n-p+1)}{(\alpha+\beta+2n-2p+2)(\alpha+\beta+2n-2p+1)},
\end{align*}
and
\begin{align*}
    a_{2p+1}&=f_{2p+1}(n)\\
    &=\frac{f_{2p-1}(n-1)}{f_{2p-1}(n)}-\frac{f_{2p}(n-1)}{f_{2p}(n)}\\
    &=\frac{-(n-p)(\alpha+n-p)}{\alpha+\beta+2n-2p+1}-\frac{-(n-p-1)(\alpha+n-p)}{\alpha+\beta+2n-2p}\\
    &=\frac{-(\alpha+n-p)(\alpha+\beta+n-p+1)}{(\alpha+\beta+2n-2p+1)(\alpha+\beta+2n-2p)}.
\end{align*}
\end{proof}


We now complete the proof of Theorem \ref{main_thm}.

\begin{proof}[Proof of Theorem \ref{main_thm}]
By Theorem \ref{elis_version}, we have
\[
\Delta_{t} = \det (s_{i+j})_{0 \leq i,j \leq t-1} = {s_0}^t(a_1a_2)^{t-1}(a_3a_4)^{t-2}\cdots(a_{2n-3}a_{2n-2}),
\]
where the $s_{k}$ are the $k$th-power-sums of the roots of $\mathscr{J}_n^{(\alpha,\beta)}(x)$ and the $a_i$ are given explicitly in the statement of the theorem.  It then follows by routine algebraic manipulation that 
\[
{s_0}^t(a_1a_2)^{t-1}(a_3a_4)^{t-2}\cdots(a_{2n-3}a_{2n-2}) = A_{t}B_{t}/C_{t},
\]
where $A_{t}$, $B_{t}$, and $C_{t}$ are given by Equations (\ref{ADEF}), (\ref{BDEF}), and (\ref{CDEF}), respectively.  
\end{proof}

\section{Hasse-Witt Invariant} \label{HW_section}


\begin{thebibliography}{99}
\bibitem[feit]{feit} W.~Feit. $\widetilde{A}_{5}$ and $\widetilde{A}_{7}$ are Galois groups over number fields. J. Algebra \textbf{104} (1986), no. 2, 231-260.
\bibitem[krattenhaler]{krattenhaler} C.~Krattenthaler. Advanced determinant calculus: a complement. Linear Algebra Appl. \textbf{411} (2005), 68-166.
\end{thebibliography}
\end{document}
